\documentclass[11pt,a4paper,twocolumn]{article}

% === Core Packages ===
\usepackage[margin=2cm]{geometry}
\usepackage[utf8]{inputenc}
\usepackage[T1]{fontenc}
\usepackage{times}
\usepackage{amsmath,amssymb}
\usepackage{graphicx}
\usepackage{booktabs}
\usepackage{float}
\usepackage{enumitem}
\usepackage{tikz}
\usepackage{pgfplots}
\pgfplotsset{compat=1.18}
\usetikzlibrary{arrows.meta,shapes.geometric,positioning,calc,decorations.markings,fit,backgrounds}

% === Citation (IEEE style) ===
\usepackage[numbers,sort&compress]{natbib}
\bibliographystyle{IEEEtranN}

% === Hyperref ===
\usepackage[hidelinks]{hyperref}
\usepackage{xurl}

% === Settings ===
\setlength{\parindent}{1em}
\setlength{\parskip}{0.3em}

% === PDF metadata ===
\hypersetup{
  pdfauthor={Pavel Kudrna},
  pdfsubject={Personix — a decentralized reputation network},
  pdfkeywords={Personix, decentralized reputation, reputation social network,
    decentralized identifiers, DID, meritocracy, censorship resistance},
}

% === Title ===
\title{\textbf{רשת מוניטין מבוזרת:\\מסגרת לארגון חברתי טבעי}}
\author{Pavel Kudrna\\Personix Project --- \url{https://personix.org}\\Prague, Czechia\\טיוטה גרסה 1 --- מרץ 2026}
\date{}

\begin{document}
\maketitle

% ============================================================
% ABSTRACT
% ============================================================
\begin{abstract}
תחושת הצדק---ההכרה שעוולות מתוקנות ושראוי נגמל---היא הכוח המלכד החזק ביותר של החברה. כאשר מוסדות ממשל מרוכזים אינם מסוגלים לספק תחושה זו משום שעלות אכיפת זכות עולה על ערכה של הזכות עצמה, הלכידות החברתית נשחקת. מבנים מוסדיים נוכחיים נכשלים ביישוב מחלקה משמעותית של סכסוכים באופן שיטתי, בדיוק כפי שתכנון מרכזי נכשל בהקצאת משאבים: דרך אסימטריית מידע, לולאות משוב חסרות וניתוק מקבלי ההחלטות מתוצאות החלטותיהם. מאמר זה מציג רשת מוניטין חברתית (RSN): שכבת תקשורת מבוזרת, בלתי ניתנת לשחיתות ועמידה בפני צנזורה הבנויה על מזהים מבוזרים (DIDs), המאפשרת למשתתפים בדויים לפרסם, לאמת ולצרוך מידע מוניטין על התנהגות בעולם האמיתי. המנגנון המרכזי נשען על שלושה עקרונות תכנון: (1)~מבנה עלות אסימטרי שבו קריאת נתוני מוניטין זולה אך כתיבה דורשת מאמץ בר-אימות; (2)~פרוטוקול בחירת מאמתים בלתי דטרמיניסטי המבוסס על גיבוב עקבי המתמחר טענות רדיקליות דרך עלויות איטרציה מסלימות; ו-(3)~מודל רשות מוניטין מונע-שוק שבו מאמתים פרטיים מהמרים על צבירת המוניטין שלהם על דיוקן של טענות שהם מאשרים. הארכיטקטורה המתקבלת מייצרת מריטוקרטיה מתהווה ללא תיאום מרכזי ומאפשרת מסלול הגירה הדרגתי והפיך ממערכות קיימות המנוהלות בידי המדינה.
\end{abstract}

% ============================================================
% 1. INTRODUCTION
% ============================================================
\section{מבוא}

\subsection{בעיית האמון המרוכז}

הממשל המודרני נשען על הנחת יסוד: שמוסדות מרוכזים יכולים לתווך אמון בין זרים בקנה מידה גדול. בתי משפט מכריעים בסכסוכים. פנקסים מאשרים בעלות. גופי רגולציה אוכפים תקנים. מוסדות אלה תוכננו בעידן שבו מידע היה נדיר, התקשורת הייתה איטית, והאצלת סמכות לגוף מרכזי הייתה מנגנון התיאום היחיד הישים. אולם ממשל מרוכז נכשל מאותן סיבות מבניות כמו תכנון מרכזי: אסימטריית מידע, לולאות משוב חסרות, וניתוק מקבלי ההחלטות מתוצאות החלטותיהם.

ההנחה נכשלת, למשל, כאשר עלות התיווך המוסדי עולה על ערך הסכסוך. שקלו דייר שמגלה כי הדירה ששכר נעדרת תעודת בטיחות חשמלית הנדרשת בחוק---מצב המהווה סכנה מיידית לחיים. זכותו החוקית של הדייר לבטל את החוזה חד-משמעית. אך עלות אכיפת אותה זכות דרך מערכת המשפט עולה על הערך הכספי של הפיקדון והנזקים המגיעים, אף בהתחשב בפיצוי סטטוטורי פוטנציאלי~\cite{czcourtfees}. התגובה הרציונלית היא לספוג את ההפסד ולפרוש.

זה אינו מקרה קצה פתולוגי. זו חוויית ברירת המחדל עבור מחלקה משמעותית של סכסוכים שבהם הנזק אמיתי אך הסיכון הכספי נופל מתחת לסף שבו אכיפה מוסדית נעשית רציונלית מבחינה כלכלית. התוצאה היא מערכת המעדיפה מבנית שחקנים רעים: מי שפוגעים באחרים בהיקפים מתחת לסף זה יכולים לפעול ללא עונש, משום שעלות חיפוש הצדק נופלת על הצד הנפגע.

הדוגמה לעיל שאובה מסביבתו המשפטית של המחבר---המשפט האזרחי הצ'כי. במסורות משפטיות אחרות---משפט מקובל מבוסס-תקדים, משפט מנהגי, מערכות מעורבות---הצדק נכשל בדרכים שונות ובמידות שונות, בהתאם למאפיינים התרבותיים והפרוצדורליים של תחום השיפוט. קוראים ככל הנראה יזהו דוגמאות מקבילות מהקשרם. מה שאוניברסלי מבנית הוא הדפוס: סכסוכים שערכם נופל מתחת לסף האכיפה הרציונלית מבחינה כלכלית מסתיימים בכך שההפסד נספג בידי הצד הנפגע. ה-RSN אינה מבטיחה צדק מושלם---היא רק מפחיתה את היתרון המבני שממנו נהנים שחקנים רעים כאשר אכיפה מוסדית אינה כלכלית.

\subsection{מדוע פתרונות קיימים אינם מספיקים}

\textbf{מסגרות זהות מבוזרת (DID)}~\cite{did2022} מספקות מנגנונים קריפטוגרפיים לניהול זהות ריבונית-עצמית אך מתמקדות בעיקר באימות זהות מבלי להתייחס לשאלה הרחבה יותר של מוניטין התנהגותי.

\textbf{אסימוני נשמה (SBTs)}~\cite{weyl2022} לוכדים את התובנה שמוניטין אינו ניתן להעברה אך מסתמכים על תשתית בלוקצ'יין עם מגבלות מדרגיות ואינם מתייחסים לתמריצים הכלכליים המסדירים את הכנסת המידע. מושגים קשורים כגון אסימוני זכות (למשל, ליברלנד) חולקים פילוסופיה דומה אך נותרים כבולים למסגרות תחום שיפוט ספציפיות.

\textbf{ארגונים אוטונומיים מבוזרים (DAOs)}~\cite{buterin2014} מיישמים ממשל דרך חוזים חכמים אך משכפלים דינמיקה פלוטוקרטית שבה ההשפעה פרופורציונלית להון. הצבעה ריבועית~\cite{buterin2019} ממתנת זאת חלקית אך מגבילה אימוץ מעשי. DAOs מתמודדים גם עם \emph{בעיית האורקל} היסודית: חוזים חכמים אינם יכולים לאמת באמינות מצבים בעולם האמיתי ללא מקור חיצוני מהימן, ולבעיה זו אין פתרון כללי בתוך ארכיטקטורת בלוקצ'יין.

אף מערכת קיימת אינה משלבת ביזור, עמידות בפני צנזורה, בדיוניות, עלות כניסה בת-אימות והסכמה מתהווה.

\subsection{תרומות}

מאמר זה תורם את התרומות הבאות:
\begin{enumerate}[nosep]
\item הגדרת \textbf{רשת המוניטין החברתית (RSN)}, פרוטוקול מבוזר המרחיב את מסגרת ה-DID לתמיכה בהחלפת מוניטין דו-צדדית.
\item \textbf{פרוטוקול בחירת מאמתים בלתי דטרמיניסטי} המבוסס על גיבוב עקבי~\cite{karger1997}.
\item \textbf{עקרון הרדיקליזם היקר}, המתמחר טענות בהתאם למרחקן מהסכמת הרשת דרך שלושה ערוצי עלות מצטברים.
\item \textbf{מודל רשות ברת-מוניטין} עם תמריצים אסימטריים שבו אישור כוזב עולה יותר בצורה קטסטרופלית מאשר עמלות לגיטימיות.
\item \textbf{מסלול הגירה הדרגתי} דרך תשתית פיסקלית מקבילה.
\end{enumerate}

% ============================================================
% 2. DESIGN PRINCIPLES
% ============================================================
\section{עקרונות תכנון}

ארכיטקטורת ה-RSN מוגבלת בחמישה עקרונות תכנון בלתי ניתנים למשא ומתן.

\textbf{המשכיות והתפתחות.} המערכת מתקיימת לצד מוסדות קיימים במהלך תקופת מעבר באורך בלתי מוגדר. המעבר חייב להיות הפיך בכל שלב.

\textbf{התנדבותיות ואחריות.} ההשתתפות התנדבותית, אך ההתנדבותיות מלווה באחריות: משתתף שנוטל שליטה על פונקציה מוסדית נוטל בו-זמנית את החובות הנלוות.

\textbf{מגוון ונייטרליות תרבותית.} ההסכמה מתהווה מאינטראקציות דו-צדדיות, לא מכללים מוגדרים מראש הנכפים בידי מתכנני המערכת.

\textbf{חוסן ועמידות בפני צנזורה.} עלות צנזור הרשת חייבת לעלות על עלות הסבלת קיומה. רק טוטליטריות מלאה---במחיר פוליטי וכלכלי הרסני---יכולה לדכא את המערכת.

\textbf{הסכמה ואכיפה.} המערכת משלבת נטיות אנושיות לקטנוניות, נקמנות וסיפוק גמולי כתכונות נושאות-משקל. התנהגות פסולה אינה אסורה; היא מתומחרת.

% ============================================================
% 3. SYSTEM OVERVIEW
% ============================================================
\section{סקירת המערכת}

\subsection{רשת המוניטין החברתית}

ה-RSN היא שכבת תקשורת מבוזרת, עמית-לעמית, שבה משתתפים מחליפים מידע בר-אימות על התנהגות בעולם האמיתי. תכונותיה המגדירות הן: תשתית מבוזרת ובלתי ניתנת לצנזור; השתתפות בדויה דרך DIDs; מבנה עלות אסימטרי (קריאה זולה, כתיבה יקרה); בר-אימות קריפטוגרפי; ו-\textbf{תמיכה בתקנים}---פורמטים קריאים-מכונה למידע המופץ דרך הרשת, לשירותים המוצעים בידי רשויות (הרכבת טענות, אישור, שירות עדות), ולפורמט המדיניות הכולל כללים להערכת מוניטין הצד שכנגד ולהערכת הסיכון של אי-התאמת מדיניות (בין התנהגות מוצהרת לבפועל).

\subsection{רכיבים ארכיטקטוניים}

ה-RSN מורכבת מארבעה רכיבים עיקריים (איור~\ref{fig:architecture}):
\begin{enumerate}[nosep]
\item \textbf{שכבת DID.} זהויות משתתפים עם כללים מוצהרים, מטא-נתוני מוניטין וניתוב רשת.
\item \textbf{פרוטוקול טענות.} פורמט מובנה לפרסום קביעות על אירועים בעולם האמיתי.
\item \textbf{רשת אימות.} פרוטוקול מבוזר להקצאת מאמתים באמצעות גיבוב עקבי.
\item \textbf{שכבת צבירת מוניטין.} שירותים המייצרים תמציות מוניטין קריאות-אדם.
\end{enumerate}

\begin{figure}[ht]
\centering
\begin{tikzpicture}[
    layer/.style={draw, minimum height=0.7cm, align=center, font=\scriptsize, fill=black!8},
    arr/.style={<->, thick, gray!60}
]
\node[layer, minimum width=5cm] (agg) at (0,2.8) {\textbf{צבירה}\\פרשנות\ $\cdot$ סיכון};
\node[layer, minimum width=2.2cm] (claim) at (-1.55,1.4) {\textbf{טענות}\\הרכבה};
\node[layer, minimum width=2.2cm] (verif) at (1.55,1.4) {\textbf{אימות}\\בחירה};
\node[layer, minimum width=5cm] (did) at (0,0) {\textbf{שכבת DID}\\זהות $\cdot$ כללים $\cdot$ ציונים};

\draw[arr] (agg.south west) -- (claim.north);
\draw[arr] (agg.south east) -- (verif.north);
\draw[arr] (claim.east) -- (verif.west);
\draw[arr] (claim.south) -- (did.north west);
\draw[arr] (verif.south) -- (did.north east);
\end{tikzpicture}
\caption{ארכיטקטורת ארבע-השכבות של RSN.}
\label{fig:architecture}
\end{figure}

\subsection{תפקידי משתתפים}

עסקת אימות מערבת עד שישה תפקידי DID נבדלים (איור~\ref{fig:roles}): \textbf{מנפיק} ($DID_I$, יוצר ומגיש את הטענה), \textbf{נושא} ($DID_S$, ה-DID שהטענה עוסקת בו---עשוי לחפוף עם המנפיק), \textbf{רשות} ($DID_A$, מאשרת את הטענה תמורת עמלה; עשויה גם להציע שירותי עדות---\emph{אופציונלי}), \textbf{עד} ($DID_{W_{1..n}}$, מנטר אנונימי של יושר המאמתים דרך קודי אתגר---\emph{אופציונלי}), \textbf{מאמת} ($DID_V$, נבחר אלגוריתמית לפרסם את הטענה), וה-\textbf{RSN} (הרשת שבה מתפרסמות טענות מאומתות). כל תפקיד ניתן להאצלה ל-DID אחר (סעיף~7.4). רשות בדרך כלל מאגדת אישור עם שירותי עדות, אך שתי הפונקציות בלתי תלויות לוגית.

\begin{figure}[ht]
\centering
\begin{tikzpicture}[
    role/.style={draw, rounded corners, minimum width=1.1cm, minimum height=0.45cm, align=center, font=\scriptsize, fill=black!8},
    opt/.style={draw, rounded corners, minimum width=1.1cm, minimum height=0.45cm, align=center, font=\scriptsize, fill=black!8, dashed},
    arr/.style={-{Stealth[length=3pt]}, thick},
    oarr/.style={-{Stealth[length=3pt]}, thick, dashed, gray},
    lbl/.style={font=\tiny, fill=white, inner sep=1pt}
]
\node[role] (I) at (0,2.4) {\textbf{מנפיק}};
\node[role] (S) at (-2.5,2.4) {\textbf{נושא}};
\node[opt] (A) at (2.5,1.2) {\textbf{רשות}};
\node[opt] (W) at (-2.5,0) {\textbf{עד}};
\node[role] (V) at (2.5,-0.6) {\textbf{מאמת}};
\node[role, fill=black!5, draw=gray] (N) at (0,-1.8) {RSN};

\draw[arr] (I) -- node[lbl] {אודות} (S);
\draw[oarr] (I) -- node[lbl,sloped] {מאשרת} (A);
\draw[oarr] (I) -- node[lbl,sloped] {אתגר} (W);
\draw[arr] (I) -- node[lbl,sloped] {בקשה} (V);
\draw[oarr] (W) -- node[lbl] {מנטר} (V);
\draw[arr] (V) -- node[lbl] {מפרסם} (N);
\end{tikzpicture}
\caption{תפקידי DID בעסקת אימות. מקווקו = אופציונלי (רשות, עד).}
\label{fig:roles}
\end{figure}

\subsection{אי-שחיתות: ארבעה כוחות}

אי-השחיתות של ה-RSN אינה נובעת ממנגנון יחיד אלא מארבעה כוחות מקבילים. אף אחד אינו מספיק לבדו; רק שילובם הופך ניסיון שחיתות ללא-רציונלי כלכלית.

\textbf{מטרה בלתי צפויה.} המאמת לכל טענה נבחר באופן בלתי דטרמיניסטי דרך גיבוב עקבי (סעיף~5.3). תוקף אינו יכול למקד מראש מאמת ספציפי לשוחד, משום שזהות המאמת היא פונקציה של תוכן הטענה ושל איטרציות קודמות, לא של בחירת המנפיק.

\textbf{מנוף מוניטין---עלייה איטית, נפילה מהירה.} מוניטין ניתן לצבור רק בהדרגה, דרך התנהגות עקבית ומתמשכת. אולם הוא יכול לאבד באופן קטסטרופלי באישור כוזב יחיד (סעיף~6.2). אסימטריה זו משמעה ששנים של מוניטין שנצבר יכולות להיהרס באישור לא-ישר יחיד; האסטרטגיה האופטימלית נותרת לדחות טענות חשודות.

\textbf{הבטחה פומבית.} כל מחזיק DID מצהיר פומבית על מדיניותו---מערך כללים קריא-מכונה שהוא מתחייב לפעול לפיו. סטייה בין ההצהרה להתנהגות בפועל ניתנת לזיהוי ומתומחרת כעבירת מוניטין חמורה יותר מההפרה הבסיסית (סעיף~7.3). צביעות לפיכך יקרה יותר מחוסר-יושר ישיר.

\textbf{אסימטריית עלות-מתקפה של הרשת.} עלות צנזור הרשת או ערעורה גדלה באופן לא-לינארי עם גודל הרשת וצפיפותה, בעוד עלות הסבלת קיומה נותרת קבועה. כדי שצנזורה תשתלם, שחקן מרוכז היה צריך להחילה על פני הרשת כולה---בדרישת אמצעים בלתי פרופורציונליים לכל תועלת נתפסת (סעיף~10).

% ============================================================
% 4. DECENTRALIZED IDENTITY MODEL
% ============================================================
\section{מודל זהות מבוזרת}

\subsection{DID עם תכונות מיוחדות}

ה-RSN מרחיבה את מפרט W3C DID Core~\cite{did2022} עם: \textbf{כללים מוצהרים} (התחייבויות התנהגותיות קריאות-מכונה), \textbf{מטא-נתוני מוניטין} (ציון התנהגותי מצטבר), ו-\textbf{ניתוב רשת} (שער בצל בר-שינוי הנפרד ממיקום הטבעת היציב).

\subsection{מחזור חיי הטענה}

טענה עוברת דרך שישה שלבים (איור~\ref{fig:lifecycle}): הרכבה, אישור רשות אופציונלי, בחירת מאמת דרך גיבוב עקבי, החלטת אימות (קבלה או דחייה עם איטרציה), פרסום, ועדכון מוניטין לכל הצדדים.

\begin{figure}[ht]
\centering
\begin{tikzpicture}[
    stp/.style={draw, rounded corners=2pt, minimum width=1.3cm, minimum height=0.45cm, align=center, font=\tiny, fill=black!5},
    arr/.style={-{Stealth[length=3pt]}, thick},
    lbl/.style={font=\tiny, fill=white, inner sep=1pt}
]
\node[stp] (comp) at (0,0) {הרכבת\\טענה};
\node[stp, dashed] (auth) at (2.0,0) {אישור\\רשות};
\node[stp] (hash) at (4.0,0) {גיבוב \&\\מיפוי טבעת};
\node[stp] (ver) at (2.0,-1.6) {בדיקת\\מאמת};
\node[stp, double] (pub) at (0,-1.6) {פורסם};
\node[stp] (rej) at (4.0,-1.6) {נדחה\\עלות $\uparrow$};

\draw[arr] (comp) -- (auth);
\draw[arr] (auth) -- (hash);
\draw[arr] (hash) -- (ver);
\draw[arr] (ver) -- node[lbl] {\checkmark} (pub);
\draw[arr] (ver) -- node[lbl] {$\times$} (rej);
\draw[arr] (rej) -- ++(0,0.8) -| (hash);
\end{tikzpicture}
\caption{מחזור חיי הטענה עם לולאת איטרציה.}
\label{fig:lifecycle}
\end{figure}

\subsection{היחס ל-W3C DID Core}

מודל הזהות של ה-RSN הוא תת-קבוצה עליונה תקינה של מפרט W3C DID Core~\cite{did2022}. כל ה-DIDs של ה-RSN הם DIDs תקפים של W3C. ההרחבות הספציפיות ל-RSN מקודדות כתכונות מסמך DID המוגדרות במפרט שיטת DID ייעודי, תואמות לתשתית קיימת כגון ION~\cite{buchner2021} ו-Ceramic~\cite{ceramic2023}.

% ============================================================
% 5. VERIFICATION AND PROPAGATION
% ============================================================
\section{אימות והפצת מידע}

\subsection{עלות כניסת המידע}

העיקרון הכלכלי המרכזי של ה-RSN הוא האסימטריה בין קריאה לכתיבה. קריאת נתוני מוניטין מתקרבת לעלות שולית אפסית עבור משתתפים החלק מרשת ה-DID ובעלי גישה התואמת את מוניטינם---חדשים חייבים להרוויח בהדרגה גישה רחבה יותר (ראו אנטי-טפילות). כתיבה---הנתפסת כמימוש בר-קיימא של מוניטין לרשת, לא כפעולה טכנית חד-פעמית---דורשת השקעה מצטברת בת-אימות על פני שלושה ממדים: \textbf{זמן} (שנות חיים המוקדשות להתנהגות עקבית, התחייבויות שמומשו וקשרים שטופחו), \textbf{אנרגיה} (מאמץ המושקע ביושר עסקי, טיפוח שותפויות ואמינות ארוכת-טווח), ו-\textbf{כסף} (השקעה בשם עצמו---פיתוח מקצועי, איכות עבודתו, השבה תמורת נזקים שנגרמו). מוניטין אינו ניתן לרכישה מיידית---הוא תוצאה של צבירה לאורך חיים שהמערכת רק הופכת לנראית וניתנת להעברה בין הקשרים. העלויות הטכניות החד-פעמיות של הפרוטוקול (חתימות קריפטוגרפיות, עמלות מאמתים) זניחות בהשוואה להשקעה בסיסית זו.

\subsection{גרף חברתי ועומק קשר}

RSN היא ביסודה רשת חברתית: כל DID מוסיף אנשי קשר (DIDs אחרים) המתירים את החיבור. לאנשי קשר אלה יש אנשי קשר משלהם, וכן הלאה. חיפוש המאמת האלגוריתמי פועל בתוך עומק מוגדר $h$ של אנשי קשר מקושרים (למשל, $h=3$: אנשי הקשר הישירים, אנשי הקשר שלהם, ואנשי קשר של אנשי קשר). ארכיטקטורה זו אינה דורשת בלוקצ'יין גלובלי יחיד---היא מעדיפה באופן טבעי את התהוותן של קהילות/אשכולות עם חפיפות לקהילות אחרות. לכל משתתף DID חייבת להיות \textbf{מדיניות} מפורסמת---מערך כללים קריא-מכונה המסדיר כיצד מוערכות בקשות נכנסות. הפרות מדיניות נרשמות ברשת כארטיפקטים שליליים.

\subsection{בחירת מאמת אלגוריתמית}

פרוטוקול האימות מעסיק גיבוב עקבי~\cite{karger1997} (איור~\ref{fig:verifier}). אימות הוא שירות בתשלום: מאמתים פועלים תמורת עמלה, ויוצרים שוק שבו התחרות מניעה תמחור, מהירות ואמינות.

\textbf{שלב 1.} מסמך הטענה משורשר עם תוצאת הגיבוב של האיטרציה הקודמת (או $\varnothing$ באיטרציה הראשונה) ומגובב:
\begin{equation}
\text{pos} = f_h(\text{claim} \| \text{prev\_hash})
\label{eq:hash}
\end{equation}

האלגוריתם חייב לכלול את \emph{המסלול המלא} של כל הבקשות והתגובות הקודמות---לא רק גיבוב של האיטרציה הקודמת. תגובתו המלאה של כל מאמת (קבלה, דחייה, מחיר מצוטט, סיבה) מצורפת למסמך הגדל, בר-אימות דרך חתימות קריפטוגרפיות בכל צעד.

\textbf{שלב 2.} הגיבוב ממופה ל-$N$ מיקומים על טבעת הגיבוב העקבי, מפוזרים באחידות (למשל, עבור $N=4$: $+0^{\circ}, +90^{\circ}, +180^{\circ}, +270^{\circ}$). מכל מיקום, חיפוש בכיוון השעון בוחר את ה-DID הקרוב ביותר כמועמד מאמת (איור~\ref{fig:multipoint}). $N$ הוא פרמטר משתנה שעשוי להיות תלוי באחוז ה-DIDs הפעילים כעת, בשעת היום, או בהיענות רשת היסטורית.

\textbf{שלב 3.} המאמת מקבל (מפרסם, ממהר על מוניטין) או דוחה (חוזר לשלב~1 עם גיבוב הדחייה). כאשר צומת אינו מגיב למרות הצהרת מצב מקוון, הבקשה הבאה חייבת לכלול את כל התגובות מכל $N$ מועמדי המאמת הקודמים.

\textbf{אנטי-טפילות.} DIDs יעד עשויים לקבוע עמלות או הגבלות גישה לשאילתות מ-DIDs חדשים בעלי מוניטין מועט. מנגנון מדורג זה (לא בינארי) מאלץ חדשים לבנות מוניטין דרך פעילות אמיתית לפני צריכה חופשית של נתוני אחרים.

\begin{figure}[ht]
\centering
\begin{tikzpicture}[
    box/.style={draw, rounded corners=2pt, minimum width=1.3cm, minimum height=0.45cm, align=center, font=\scriptsize, fill=black!5},
    dec/.style={draw, diamond, aspect=2.2, minimum width=0.9cm, align=center, font=\scriptsize, fill=black!5},
    arr/.style={-{Stealth[length=3pt]}, thick},
    lbl/.style={font=\tiny, fill=white, inner sep=1pt}
]
\node[box] (iss) at (0,0) {מנפיק};
\node[box] (hash) at (0,-1.1) {$f_h$(claim $\|$\\prev\_hash)};
\node[box] (ring) at (0,-2.2) {מיפוי טבעת\\בכיוון השעון};
\node[box, dashed] (spam) at (0,-3.2) {פיקדון\\אנטי-ספאם};
\node[dec] (check) at (0,-4.4) {בדיקת\\מדיניות};
\node[box] (rej) at (2,-5.6) {איטרציה\\הבאה};
\node[box] (fee) at (0,-5.6) {עמלה סופית =\\$f$(נתונים, מוניטין, מדיניות)};
\node[box, double] (pub) at (0,-6.8) {פורסם};

\draw[arr] (iss) -- (hash);
\draw[arr] (hash) -- (ring);
\draw[arr] (ring) -- (spam);
\draw[arr] (spam) -- (check);
\draw[arr] (check) -- node[lbl] {$\times$} (rej);
\draw[arr] (check) -- node[lbl] {\checkmark} (fee);
\draw[arr] (fee) -- (pub);
\draw[arr] (rej.east) -- ++(0.5,0) |- (hash.east);
\end{tikzpicture}
\caption{פרוטוקול בחירת מאמת. פיקדון האנטי-ספאם אופציונלי ועשוי להיות בר-החזר. העמלה הסופית משולמת רק בקבלה.}
\label{fig:verifier}
\end{figure}

כל דחייה מצרפת את תגובתו המלאה של המאמת (כולל סיבות ותנאים) למסמך הטענה, ומרחיבה את תוכנו. מהמסמך המורחב, גיבוב חדש מחושב באופן בלתי דטרמיניסטי, ממופה למיקום טבעת שונה ובוחר מאמת שונה. תיעוד המסלול המלא הוא חובה---המנפיק אינו יכול לחזות או להשפיע על המאמת הבא. אם מאמת מתעלם מבקשה למרות מדיניות מוצהרת תואמת, עדים (אם קיימים) רושמים זאת, והמנפיק עשוי לצרף ראיות עדות בעת הפרסום הסופי.

\begin{figure}[ht]
\centering
\begin{tikzpicture}[scale=0.65]
% Ring
\draw[thick] (0,0) circle (2.2cm);
% DID nodes on ring
\foreach \a in {10,55,80,120,150,195,230,260,305,340} {
  \fill[black!40] (\a:2.2) circle (1.5pt);
}
% Hash offset arrow pointing to initial position
\draw[-{Stealth[length=3pt]}, thick, black!60] (0,0) -- node[font=\tiny,above,sloped,pos=0.6] {$f_h$} (37:2.2);
\fill[black] (37:2.2) circle (2pt);
\node[font=\tiny,anchor=south west] at (37:2.35) {$h_0$};
% N=4 points offset from h0 by +0, +90, +180, +270
\foreach \offset/\l in {0/P1,90/P2,180/P3,270/P4} {
  \pgfmathsetmacro{\ang}{37+\offset}
  \node[fill=black, circle, inner sep=1.5pt] (p\offset) at (\ang:2.2) {};
  \node[font=\tiny,font=\bfseries] at (\ang:2.75) {\l};
  % clockwise search arc
  \draw[-{Stealth[length=2.5pt]}, thick, gray] (\ang:2.2) arc (\ang:\ang+30:2.2);
}
\node[font=\scriptsize, align=center] at (0,0) {טבעת\\גיבוב};
\node[font=\tiny, align=center] at (0,-3.2) {$h_0 = f_h(\text{claim})$, ואז $+0^{\circ},+90^{\circ},+180^{\circ},+270^{\circ}$\\כל נקודה $\rightarrow$ חיפוש בכיוון השעון};
\end{tikzpicture}
\caption{בחירה רב-נקודתית ($N{=}4$). הגיבוב $h_0$ קובע את ההיסט; 4 נקודות מפוזרות באחידות מ-$h_0$.}
\label{fig:multipoint}
\end{figure}

\subsection{פרוטוקול העדים}

תפקיד ה-\textbf{עד} מספק אחריותיות אנונימית ליושר המאמתים דרך פרוטוקול קוד-אתגר קריפטוגרפי:

\textbf{שלב W1.} המבקש חותם על בקשת האימות ושולח אותה ל-$N_w$ עדים---אנשי קשר מהרשת החברתית של המבקש, או ספקי שירות עדות מתמחים הפועלים תמורת עמלה.

\textbf{שלב W2.} כל עד $w_i$ חותם על הבקשה החתומה כולה, שומר את החתימה המקורית $\sigma_i$, ומחשב קוד אתגר $c_i = h(\sigma_i)$. קוד האתגר מצורף לבקשה.

\textbf{שלב W3.} הבקשה, הנושאת כעת $N_w$ קודי אתגר, נשלחת למאמת. המאמת רואה את הקודים אך \emph{אינו יכול לקבוע} מי העדים או האם הקודים תואמים לחתימות אמיתיות.

\textbf{שלב W4 (מאמת ישר).} אם המאמת מגיב בהתאם למדיניותו המוצהרת, קודי האתגר נותרים אטומים. עדים לעולם אינם נחשפים. לא מתפרסמים נתונים נוספים.

\textbf{שלב W5 (הפרת מדיניות).} אם המאמת מפר את מדיניותו (מתעלם מהבקשה, מגיב בחוסר עקביות), המבקש מחזיק את החתימות המקוריות $\sigma_i$ של העדים. אלה ניתנות לפרסום כנתוני לוואי: כל אחד יכול לאמת $c_i = h(\sigma_i)$, ובכך להוכיח שהעד היה נוכח. המבקש יכול לפרסם באופן עצמאי---המאמת כבר חתם על קודי האתגר בתגובתו.

\textbf{תכונת הבלוף.} המבקש עשוי להחליף חלק מקודי האתגר או את כולם בערכים אקראיים. המאמת אינו יכול להבחין בין קודים אמיתיים (מגובים ברשויות בעלות מוניטין גבוה) לרעש. \emph{אי-ודאות} זו היא מנגנון האכיפה: כל בקשה נושאת את הסיכון שרשות מכובדת צופה באנונימיות. עלות שמירת לחץ זה קרובה לאפס (יצירת מספרים אקראיים), בעוד העלות הפוטנציאלית של חוסר-יושר עבור המאמת קטסטרופלית. התנהגות ישרה מתומרצת אף בהיעדר עדים בפועל.

\textbf{רשות כעד.} רשות המאשרת טענה עשויה לאגד שירותי עדות: ``אני מאשרת את טענתך ומשמשת עד אנונימי במהלך האימות.'' זהו מודל עסקי טבעי---לרשות כבר יש הקשר. אולם שני התפקידים בלתי תלויים לוגית.

\subsection{עקרון הרדיקליזם היקר}

שלושה ערוצי עלות מצטברים בו-זמנית (איור~\ref{fig:radicalism}):

\textbf{ערוץ 1: עלות איטרציה.} כל דחייה מאלצת איטרציה חדשה הצורכת זמן ומשאבים. צמיחה לינארית.

\textbf{ערוץ 2: התנפחות מסמך.} כל איטרציה מוסיפה את תגובת המאמת המלאה למסמך הטענה. הצמיחה לינארית, אך עם היסט בסיס: המטען הראשוני (תוכן הטענה, אישור רשות, חתימות קריפטוגרפיות) כבר בעל גודל לא-טריוויאלי $B_0$. גודל כולל לאחר $k$ איטרציות: $S(k) = B_0 + k \cdot \Delta$, כאשר $\Delta$ הוא גודל התגובה הממוצע.

\textbf{ערוץ 3: פרמיית סיכון המאמת.} הסיכון סובייקטיבי. המאמת מעריך האם המדיניות המוצהרת של ה-DID המבקש מתנגשת עם האינטרסים המסחריים, החברתיים או הפוליטיים של המאמת עצמו. הפרמיה עשויה להסלים באופן מעריכי עם המרחק הנתפס מ``האידיאל'' של המאמת: $P(d) = \alpha \cdot e^{\beta d}$, כאשר $d$ הוא מדד המרחק הסובייקטיבי של המאמת. מעבר לגבול איסור אישי, המאמת עשוי לסרב לחלוטין.

\begin{figure}[ht]
\centering
\begin{tikzpicture}
\begin{axis}[
    width=\columnwidth,
    height=4.5cm,
    xlabel={\footnotesize רדיקליזם},
    ylabel={\footnotesize עלות יחסית (לוג)},
    ymode=log,
    xmin=0, xmax=100,
    ymin=1, ymax=10000,
    grid=major,
    grid style={gray!20},
    legend style={at={(0.03,0.97)},anchor=north west,font=\tiny,draw=gray!50},
    tick label style={font=\tiny},
    label style={font=\footnotesize},
]
\addplot[black, thick, domain=0:100, samples=50] {1 + 0.3*x};
\addlegendentry{עלות איטרציה (לינארית)}
\addplot[black, thick, dashed, domain=0:100, samples=50] {1 + 0.15*x};
\addlegendentry{התנפחות מסמך (לינארית)}
\addplot[black, thick, dotted, line width=1.2pt, domain=0:100, samples=100] {exp(0.08*x)};
\addlegendentry{פרמיית סיכון (מעריכית)}
\end{axis}
\end{tikzpicture}
\caption{הסלמת עלות על פני שלושה ערוצים. פרמיית הסיכון שולטת ברדיקליזם גבוה.}
\label{fig:radicalism}
\end{figure}

באופן קריטי, זו אינה צנזורה. אף טענה אינה אסורה. המערכת מתמחרת סיכון (טבלה~\ref{tab:costs}).

\begin{table}[ht]
\centering
\caption{השוואת עלות על פני סוגי טענות.}
\label{tab:costs}
\footnotesize
\begin{tabular}{@{}lcccc@{}}
\toprule
\textbf{סוג} & \textbf{איטר.} & \textbf{מסמך} & \textbf{עמלה} & \textbf{סה"כ} \\
\midrule
אמין & $k_{\min}$ & $B_0$ & $P_{\min}$ & נמוך \\
שנוי במחלוקת & $k_{\text{med}}$ & $B_0{+}k\Delta$ & $\alpha e^{\beta d}$ & בינוני \\
רדיקלי & $k_{\max}$ & $\gg B_0$ & $\gg P_{\min}$ & גבוה מאוד \\
מרמתי & $\to\infty$ & --- & --- & בלתי ניתן לפרסום \\
\bottomrule
\end{tabular}
\end{table}

% ============================================================
% 6. REPUTATION AND RISK
% ============================================================
\section{מוניטין והערכת סיכון}

\subsection{מודל צבירת המוניטין}

מוניטין מציג שלוש תכונות: \textbf{צבירה איטית} (צמיחה הדרגתית דרך התנהגות עקבית), \textbf{הרס מהיר} (מרמה יחידה יכולה להפחית את הציון באופן קטסטרופלי), ו-\textbf{הערכה פרטנית} (כל צד צורך עשוי להחיל פונקציית צבירה משלו).

\subsection{רשויות ברות-מוניטין}

רשות ברת-מוניטין בוחנת טענות, ואם מרוצה, חותמת עליהן במשותף---ממהרת על מוניטינה שלה (איור~\ref{fig:authority}). האסימטריה מתוכננת: צבירת מוניטין איטית (הדרגתית $+\delta$ לכל אישור ישר), בעוד אובדנו קטסטרופלי ($-\Delta \gg \delta$ לכל אישור כוזב; להמחשה, למשל, $+2\%$ לעומת\ $-70\%$). האסטרטגיה האופטימלית היא תמיד לדחות טענות חשודות. הערכים הספציפיים של $\delta$ ו-$\Delta$ הם פרמטרים של המערכת.

\begin{figure}[ht]
\centering
\begin{tikzpicture}[
    box/.style={draw, rounded corners=2pt, minimum width=2cm, minimum height=0.6cm, align=center, font=\scriptsize, fill=black!5},
    arr/.style={-{Stealth[length=4pt]}, thick}
]
\node[box] (exam) at (0,0) {הרשות בוחנת\\מוניטין: $R$};
\node[box] (true) at (-2,-1.3) {אמת: $R{\to}R{+}\delta$\\עוד לקוחות};
\node[box] (false) at (2,-1.3) {שקר: $R{\to}R{-}\Delta$\\לקוחות עוזבים};
\node[box] (insight) at (0,-2.6) {רווח: איטי ($+\delta$)\\אובדן: מיידי ($-\Delta$)};

\draw[arr] (exam) -- node[left,font=\tiny] {אמת} (true);
\draw[arr] (exam) -- node[right,font=\tiny] {שקר} (false);
\draw[arr] (true) -- (insight);
\draw[arr] (false) -- (insight);
\end{tikzpicture}
\caption{רשות ברת-מוניטין---תמריצים אסימטריים.}
\label{fig:authority}
\end{figure}

\subsection{מוניטין רב-ממדי}

מוניטין אינו סקלר אלא וקטור על פני ממדים מרובים: אמינות פיננסית, יושרה אישית, כשירות מקצועית, מעורבות אזרחית ואחרים (איור~\ref{fig:reputation-radar}). ספקי הערכה שונים מטילים וקטור זה באופן שונה בהתאם להקשר---מלווה מעדיף אמינות פיננסית, מעסיק כשירות מקצועית, שכן התנהגות אזרחית. אין ציון מוניטין ``נכון'' יחיד; רק נקודות מבט.

\begin{figure}[ht]
\centering
\begin{tikzpicture}[scale=0.55]
\foreach \a/\l in {90/פיננסי,162/יושרה,234/מקצועי,306/אזרחי,18/חברתי} {
  \draw[gray!40] (0,0) -- (\a:2.5);
  \node[font=\tiny, align=center] at (\a:3.1) {\l};
  \foreach \r in {0.8,1.6,2.4} {
    \fill[black!15] (\a:\r) circle (0.5pt);
  }
}
\draw[thick, fill=black!10] (90:2.0) -- (162:1.2) -- (234:1.8) -- (306:0.9) -- (18:1.5) -- cycle;
\draw[thick, dashed] (90:1.4) -- (162:2.1) -- (234:0.8) -- (306:2.0) -- (18:1.1) -- cycle;
\node[font=\tiny] at (0,-3.3) {מלא: DID A \quad מקווקו: DID B};
\end{tikzpicture}
\caption{וקטורי מוניטין רב-ממדיים. DIDs שונים מציגים פרופילים שונים על פני הממדים.}
\label{fig:reputation-radar}
\end{figure}

\subsection{הערכת סיכון מודמוקרטת}

ה-RSN מדמוקרטת הערכת סיכון שיטתית---יכולת המרוכזת כיום במוסדות פיננסיים אך ישימה הרבה מעבר לבנקאות. תאגידי ענק תעשייתיים המעריכים אמינות ספקים, חברות ביטוח המתמחרות סיכון התנהגותי, בדיקת נאותות למיזוגים ורכישות, הערכת אורך חיי ציוד בייצור---בכל מקום שבו סיכון הצד שכנגד דורש הערכה, ה-RSN מספקת חלופה מבוזרת ומונעת-שוק. שוק של שירותי פרשנות מתרגם נתוני מוניטין רב-ממדיים גולמיים לתמציות ניתנות-לפעולה, ומנגיש הערכת צד שכנגד לכל משתתף בעלות שולית.

% ============================================================
% 7. CONSENSUS AND DISPUTE RESOLUTION
% ============================================================
\section{הסכמה ויישוב סכסוכים}

\subsection{מודל צדק מבוזר}

ה-RSN ממסגרת מחדש את הצדק כבעיית משא ומתן דו-צדדית. איום נזק מוניטין קבוע ובר-אימות הוא מרתיע יעיל יותר מהליכים משפטיים שהצד הנפגע אינו יכול להרשות לעצמו.

\subsection{הסכמה מתהווה}

כל משתתף מקיים סף שביעות רצון אישי. ההסכמה המצרפית של הרשת מתהווה כממוצע הסטטיסטי של אלפי משאים ומתנים דו-צדדיים (איור~\ref{fig:consensus}). תגובה הרבה מעל ההסכמה (ברברית) יקרה לפרסום. תגובה סמוך להסכמה (פרופורציונלית) זולה. ההסכמה אינה כלל---היא אות מחיר.

\begin{figure}[ht]
\centering
\begin{tikzpicture}[
    person/.style={draw, rounded corners=2pt, minimum width=1.5cm, minimum height=0.5cm, align=center, font=\scriptsize, fill=black!5},
    arr/.style={-{Stealth[length=4pt]}, thick}
]
\node[person] (a) at (-2,0) {קפדני\\(גבוה)};
\node[person] (b) at (0,0) {פרגמטי\\(בינוני)};
\node[person] (c) at (2,0) {סלחני\\(נמוך)};
\node[person, fill=black!10, minimum width=3cm] (cons) at (0,-1.2) {הסכמת הרשת\\= ממוצע סטט.};
\node[person] (exp) at (-2,-2.4) {ברברי\\יקר};
\node[person, double] (prop) at (0,-2.4) {פרופורציונלי\\זול};
\node[person] (neg) at (2,-2.4) {רשלני\\יקר};

\draw[arr] (a) -- (cons);
\draw[arr] (b) -- (cons);
\draw[arr] (c) -- (cons);
\draw[arr] (cons) -- (exp);
\draw[arr] (cons) -- (prop);
\draw[arr] (cons) -- (neg);
\end{tikzpicture}
\caption{הסכמה מתהווה מספי שביעות רצון פרטניים.}
\label{fig:consensus}
\end{figure}

\subsection{כיול ענישה}

כל מחזיק DID מצהיר פומבית על תגובות לקטגוריות התנהגות פסולה ספציפיות. הצהרות קשות או מקלות באופן בלתי פרופורציונלי מושכות אותות מוניטין שליליים. הצהרות פרופורציונליות מתקבלות ללא חיכוך---מערכת ענישה מכיילת-עצמית.

הצהרות הסכמה כפופות בעצמן לזיהוי צביעות: התחייבות הצהרתית לעמדת הסכמה תוך התנהגות לא-עקבית מהווה עבירת מוניטין חמורה יותר מהסטייה הבסיסית, שכן היא מדגימה הטעיה מכוונת.

\subsection{האצלה}

כל הפעילויות בתוך DID---למעט חריג אחד---ניתנות להאצלה ל-DID אחר. \textbf{תפקיד הנושא---ה-DID שהתנהגותו נושא הטענה---אינו ניתן להאצלה; הוא כבול באופן בלתי ניתן להעברה למחזיק הזהות שהטענה מתייחסת אליו.} התפקידים הפעילים האחרים---מנפיק, רשות, עד, מאמת---ניתנים להעברה ל-DID נציג ייעודי, בין אם לאימות, הגשת טענות, השתתפות בהסכמה או יישוב סכסוכים. ההאצלה ניתנת לביטול בכל עת. זה מאפשר התמחות (למשל, שירותי אימות מקצועיים) בעוד המחזיק שומר על שליטה ואחריות סופיות. ההאצלה חלה על פני המערכת כולה וחיונית למדרגיות.

\subsection{ממוסר פרטני לחוזה חברתי מתהווה}

המדיניות המוצהרת של כל DID מייצגת פורמליזציה של מוסר פרטני: מה שהמחזיק רואה כמקובל, כיצד הוא מגיב להתנהגויות ספציפיות, מה הוא דורש מהצדדים שכנגד. מדיניות זו אינה נכפית בידי מתכנן מערכת---היא נבחרת בידי כל משתתף ומובעת בצורה קריאה-מכונה.

פורמליזציה זו מאפשרת התהוות תלת-שלבית. ראשית, מוסר פרטני מקודד כ-\textbf{מדיניות}---מערך כללים מוצהרים, ספים ופונקציות תגובה שה-DID מתחייב לפעול לפיהם. שנית, המצרף של כל המדיניות על פני הרשת מייצר \textbf{אתיקה מתהווה}---נורמות ניתנות-לתצפית סטטיסטית שאף משתתף יחיד לא תכנן אך המתעוררות מהאינטראקציה של אלפי עמדות מוסריות פרטניות~\cite{durkheim1893}. שלישית, אתיקה מתהווה זו, המובעת כנורמות התנהגות משותפות הנאכפות דרך אותות מחיר דו-צדדיים, מהווה \textbf{חוזה חברתי מדיד}---לא מבנה תיאורטי שאיש לא חתם עליו~\cite{rawls1971}, אלא מערך כללים ניתן-לתצפית אמפירית המגובה בהתנהגות בפועל.

תהליך זה משקף ממצאים מתאוריית היסודות המוסריים~\cite{haidt2012}: המוסר האנושי אינטואיטיבי, פלורליסטי ומשתנה תרבותית. ה-RSN אינה דורשת הסכמה מוסרית כקלט; היא מייצרת הסכמה התנהגותית כפלט. החוזה החברתי המתקבל אינו סטטי---הוא מתפתח ככל שמשתתפים מצטרפים, עוזבים ומתאימים את מדיניותם בתגובה למשוב הרשת. בשונה מתאוריית החוזה החברתי הקלאסית, חוזה זה ניתן לביטול, ניתן לתצפית וניתן לאכיפה ללא ריבון.

% ============================================================
% 8. ECONOMIC MODEL
% ============================================================
\section{מודל כלכלי}

הסעיפים הקודמים תיארו כלי---מנגנוני האימות, מבנה העלות וההסכמה המתהווה של ה-RSN. סעיף זה מתאר מתודולוגיה ליישום כלי זה על מעבר פיסקלי וממשלי. הכלי הוא רב-תכליתי; המתודולוגיה שלהלן היא יישום אפשרי אחד.

\subsection{מבנה תמריצים}

מוניטין כהון יוצר תמריצים ישירים להתנהגות ישרה. בתפעול רגיל, משתתפים בונים מוניטין באופן טבעי דרך עסקאות יומיומיות והתחייבויות שמומשו. ההוצאה העיקרית היא על טענות \emph{שליליות}---רישום נזק שנגרם בידי אחרים. אסימטריה זו מתמרצת התנהגות מועילה ואמינה: להיות ישר אינו עולה דבר נוסף, בעוד להיות מזיק יוצר עקבה מתועדת שאחרים ישלמו כדי לשמר. צביעות---הצהרת כללים מבלי לפעול לפיהם---היא מצב הכשל היקר ביותר, שכן היא מדגימה הטעיה מכוונת.

\subsection{מערכת פיסקלית מקבילה}

שלושה רכיבים מאפשרים לחץ תחרותי: (1)~\textbf{מס פשוט}---שיעור פרופורציונלי יחיד ללא חריגים, בתחילה מעט מתחת לשיעור האפקטיבי הקיים; (2)~\textbf{רישום הוצאות אלקטרוני (ESR)}---היפוך מודל ה-EET הצ'כי כך שאזרחים מפקחים על הוצאות המדינה; (3)~\textbf{הקצאת מס בהכוונת אזרח}---החל מאחוזים בודדים בשנה הראשונה ומגיע, לאחר עשור, לכל מקום מהעשירונים הכפולים הנמוכים ועד הגירה מלאה למערכת בהכוונת אזרח, בהתאם לקצב האימוץ. הקצב בפועל תלוי במהירות שבה מתעוררות חלופות שוק חופשי לשירותי המדינה ובנכונות המדינה לשתף פעולה עם המעבר; פונקציית הזמן אינה חייבת להיות לינארית, ואין סיבה לקבוע גבול עליון לנקודה שאליה עשוי האימוץ להגיע בסופו של דבר.

% ============================================================
% 9. MIGRATION PATH
% ============================================================
\section{מסלול הגירה}

ההגירה מתקדמת דרך שלושה שלבים (איור~\ref{fig:migration}): \textbf{שלב~1: שכבת-על} (RSN כשכבת מידע, המדינה היא הספק היחיד), \textbf{שלב~2: תחרות} (RSN מספקת חלופות פונקציונליות, שימוש דו-מערכתי), ו-\textbf{שלב~3: מעבר} (RSN עולה בביצועים על מקבילות המדינה, הגירה מרצון). המעבר הפיך בכל שלב.

\begin{figure}[ht]
\centering
\begin{tikzpicture}[
    phase/.style={draw, rounded corners=3pt, minimum width=1.8cm, minimum height=0.8cm, align=center, font=\scriptsize, fill=black!5},
    arr/.style={-{Stealth[length=5pt]}, thick}
]
\node[phase] (p1) at (0,0) {\textbf{שלב 1}\\שכבת-על};
\node[phase] (p2) at (2.2,0) {\textbf{שלב 2}\\תחרות};
\node[phase] (p3) at (4.4,0) {\textbf{שלב 3}\\מעבר};

\draw[arr] (p1) -- (p2);
\draw[arr] (p2) -- (p3);
\draw[arr, dashed, gray] (p3.south) -- ++(0,-0.6) -| node[below,font=\tiny,pos=0.25] {נסיגה} (p2.south);
\end{tikzpicture}
\caption{הגירה בת-שלושה-שלבים---הפיכה בכל שלב.}
\label{fig:migration}
\end{figure}

מנגנון הלחץ העיקרי הוא פיסקלי: כאשר משלמי מס מותנים מגיעים ל``מיעוט משמעותי כלכלית $X$''---קבוצה גדולה דיה שדיכוי נגדה גורם נזק כלכלי לא-טריוויאלי ואי-ציות אזרחי הופך לאיום אמין---המדינה נכנסת למשא ומתן. להמחשה, אנו משתמשים ב-$X \approx 15\%$ מהתמ"ג, אך הסף בפועל תלוי בכלכלה הספציפית.

% ============================================================
% 10. SECURITY ANALYSIS
% ============================================================
\section{ניתוח אבטחה}

אנו שוקלים חמש קטגוריות יריב: שחקנים רעים בודדים, קבוצות מאורגנות, יריבים תאגידיים, יריבים ברמת מדינה, ויריבים ברמת הפרוטוקול.

\textbf{עמידות בפני סיביל}~\cite{douceur2002}. בניית מוניטין דורשת אינטראקציה מתמשכת ובת-אימות. עלות מתקפת סיביל מוצלחת גדלה לינארית עם זהויות מזויפות וריבועית עם רמת המוניטין המבוקשת. תפעול DIDs מרובים במקביל אפשרי אך יקר מעצם התכנון---כל זהות חייבת לצבור מוניטין באופן עצמאי דרך פעילות אמיתית. בחברות חופשיות עלות זו מרתיעה שימוש לרעה מזדמן; בדיקטטורות, DIDs מקבילים הופכים לכלי הישרדות המאפשר התנגדות מתאיית, תיאום מחתרתי וניווט בטוח יותר בשוק השחור.

\textbf{הגנת קנוניה.} דפוסי אישור הדדי בקרב קבוצות סגורות ניתנים לזיהוי דרך ניתוח גרף חברתי. פונקציות צבירת מוניטין מנכות טענות ממקורות מאושכלים בצפיפות.

\textbf{יריב ברמת מדינה.} ניתוב בצל, היעדר תשתית מרוכזת, ופיזור תפקיד המאמת על פני בסיס המשתתפים מבטיחים שדיכוי דורש אמצעים בלתי פרופורציונליים לכל תועלת.

\textbf{פרטיות מול אחריותיות.} בדיוניות---דרך אמצע בין אנונימיות לחשיפת זהות---מאפשרת ל-DIDs לצבור מוניטין משמעותי מבלי לחשוף את זהותו האמיתית של המחזיק.

% ============================================================
% 11. PRIVACY
% ============================================================
\section{שיקולי פרטיות}

מודל הפרטיות של ה-RSN בנוי על בדיוניות. המחזיק שולט בחשיפת הזהות: מינימלית (פסבדונים טהור), סלקטיבית (אישורים ספציפיים מקושרים), או מלאה (זהות חוקית משויכת). טענות שפורסמו הן קבועות---המערכת תומכת בתיקון הקשרי אך לא במחיקה. מחזיק עשוי לנטוש DID שנפרץ ולהתחיל מחדש, תוך ויתור על מוניטין שנצבר.

% ============================================================
% 12. RELATED WORK
% ============================================================
\section{עבודה קשורה}

מפרט W3C DID Core~\cite{did2022} ורשת Ceramic~\cite{ceramic2023} מספקים את יסוד הזהות. EigenTrust~\cite{kamvar2003} הדגים צבירת מוניטין מבוזרת ללא רשות מרכזית. SBTs~\cite{weyl2022} הציגו אסימונים חברתיים בלתי ניתנים להעברה. ה-RSN נבדלת בהדגשה שלה על עלות כלכלית כמסנן איכות ובמנגנון שלה לתמחור רדיקליזם.

DAOs~\cite{buterin2014} והצבעה ריבועית~\cite{buterin2019} הדגימו את היתכנותו של ממשל מבוזר. ה-RSN גוזרת הסכמה ממשאים ומתנים על מחיר דו-צדדיים במקום מהצבעה.

ההצעה נשענת על התאוריה האנרכו-קפיטליסטית~\cite{rothbard1973,hoppe2001} לאספקת שירות פרטית אך נבדלת בשני היבטים קריטיים: היא מתכננת עבור נקמנות ודחפים גמוליים במקום להניח רציונליות, והיא מציעה מעבר הדרגתי במקום מהפכה. למושג החוזים החברתיים המתהווים יש קדמונים בתאוריה של הייק על סדר ספונטני~\cite{hayek1973}.

\textbf{אנרכו-אגוריזם.} ה-RSN חולקת מכנה משותף משמעותי עם הכלכלה-הנגדית האגוריסטית~\cite{konkin2006}---בפרט ההדגשה על בניית מוסדות מקבילים והחלפה מרצון. אולם אגוריזם נוטה להניח אינטרס עצמי רציונלי כמנגנון תיאום מספיק ולעיתים קרובות חסר מסלול הגירה קונקרטי ממבנים מדינתיים קיימים. ה-RSN משלבת במפורש נטיות התנהגותיות לא-רציונליות ומספקת מנגנון לחץ פיסקלי למעבר הדרגתי.

\textbf{מריטוקרטיה.} ה-RSN עשויה לייצג תיאור פונקציונלי ראשון של האופן שבו מריטוקרטיה~\cite{young1958} יכולה להתהוות כתכונה מערכתית ולא כסיסמה. מערכות מריטוקרטיות מסורתיות נכשלות משום שהן דורשות רשות מרכזית להגדיר ולאכוף ``ראוי''. ב-RSN, הראוי מתהווה מהערכות דו-צדדיות: מי שמספקים ערך צוברים מוניטין; אף ועדה אינה מחליטה למי יש ראוי.

\textbf{קניין כפריבילגיה.} ה-RSN נבדלת ביסודה מטיפול האסכולה האנרכו-קפיטליסטית והאוסטרית בזכויות קניין כטבעיות ומוחלטות. במסגרת ה-RSN, קניין הוא \emph{פריבילגיה}---הכרה חברתית שניתן, במקרים קיצוניים, לשלול בידי הקהילה כאשר משתתף מפר ביסודו נורמות משותפות (בגידה, סירוב להגן על הקהילה במשבר קיומי, ניצול שיטתי). זו אינה החרמה מדינתית אלא שלילת הכרה חברתית בידי רשת היחסים הדו-צדדיים.

% ============================================================
% 13. LIMITATIONS
% ============================================================
\section{דיון ומגבלות}

\textbf{התנעה קרה.} ערכה של ה-RSN תלוי באפקטי רשת; מאמצים מוקדמים ניצבים בפני ערך מוגבל עד שההשתתפות מגיעה לסף. אולם, אף משלבים מוקדמים, רשת ה-DID משמשת מקור מידע משלים שימושי. הערכת מוניטין מוקדמת והערכת רשויות צפויות להתפתח בהדרגה עם תחכום גובר.

\textbf{אנטי-טפילות ובקרת גישה.} DIDs יעד עשויים להטיל עמלות מדורגות והגבלות גישה על שאילתות מ-DIDs בעלי מוניטין נמוך. זה אינו בינארי אלא סקאלה רציפה: ככל ש-DID שואל בעל פחות מוניטין, כך הוא מקבל פחות מידע, ממתין זמן רב יותר, ומשלם יותר. מנגנון זה מתמרץ השתתפות אמיתית על פני צריכה פסיבית.

\textbf{אי-שוויון בגישה.} עלות פרסום הטענות עשויה לסנן תלונות לגיטימיות ממשתתפים מוחלשים כלכלית. הפחתה: כל משתתף משמש בו-זמנית מאמת פוטנציאלי (או מאציל תפקיד זה) ומרוויח עמלות אימות. לאורך אופק זמן מספיק, משתתף לא-רדיקלי אמור להתקרב לניטרליות פיננסית---הכנסת האימות מקזזת בקירוב את עלויות הפרסום. זו ציפייה סטטיסטית, לא ערובה.

\textbf{אי-שוויון מוניטין.} מאמצים מוקדמים צוברים יתרונות שעשויים ליצור מחסומים למגיעים מאוחר. אולם, הערכת מוניטין מבוצעת בידי ספקי צד שלישי שעשויים לבחור להפחית את יתרון-הראשוניות דרך אלגוריתמי הצבירה שלהם. להיות ראשון עם מוניטין טוב הוא יתרון כלכלי השוואתי לגיטימי---וכך במתכוון.

\textbf{שאלות פתוחות.} פונקציות עלות אופטימליות, תקני צבירה, ממשל פרוטוקול, ואינטראקציה עם נתוני מוניטין סינתטיים שנוצרו בידי בינה מלאכותית נותרים תחומים לעבודה עתידית.

מאמר זה אינו טוען שה-RSN תייצר תוצאות אופטימליות בכל הנסיבות. הוא טוען שה-RSN מייצגת חלופה \emph{ישימה}---ניתנת ליישום טכני, בת-קיימא כלכלית, ותואמת חברתית לנטיות ההתנהגות האנושיות כפי שהן באמת.

% ============================================================
% 14. CONCLUSION
% ============================================================
\section{מסקנה}

מאמר זה הציג את רשת המוניטין החברתית, פרוטוקול מבוזר המאפשר החלפת מוניטין בדויה ובת-אימות. עקרון הרדיקליזם היקר מבטיח שטענות מרמתיות מתומחרות אל מחוץ לתמונה ללא צנזורה. מסלול ההגירה המוצע מאפשר מעבר הדרגתי והפיך ממוסדות קיימים. התכנון משלב את הטבע האנושי כפי שהוא---כולל קטנוניות, נקמנות והרצון לסיפוק גמולי---במקום לדרוש טרנספורמציה אידאולוגית. התוצאה אינה אוטופיה אלא מערכת שבה הצדק נגיש, המוניטין נרכש, ועלות ההתנהגות הפסולה נישאת בידי הצד שגרם לה.

% ============================================================
% REFERENCES
% ============================================================
\bibliography{references}

\end{document}
